GLev 密文是由多个 GLWE 密文组成的列表，分别加密明文列表 $\dfrac{q}{\beta^1}M, \dfrac{q}{\beta^2}M, \ldots, \dfrac{q}{\beta^l}M$，其中 $M$ 是编码为多项式的明文。GLev 密文中的第 $i$ 个 GLWE 密文使用不同的明文尺度因子 $\Delta_i = \dfrac{q}{\beta^i}$。\autoref{fig:glev} 直观展示了 GLev 密文的结构。


对于这里的整数环权重，要求整数 $\beta\ge2$ 且 $\beta^l\mid q$。只有在期望第一行恢复取自 $\{0,\ldots,t-1\}$ 的系数代表元时，才需要 $t\le\beta$。GLev 也可作为一般环消息的 gadget 密文使用；在这种用法中，并不要求每一行都能单独解码原消息。若要用降幂 gadget 精确重构每个剩余，还须满足 $q=\beta^l$，如 \autoref{sec:decomp} 所述。


\subsection{加密}
\label{subsec:glev-enc}

\begin{tcolorbox}[title={\textbf{\tboxlabel{\ref*{subsec:glev-enc}} GLev 加密}}]


$\textsf{GLev}_{S, \sigma}^{\beta, l}(M) = \Bigl \{\textsf{GLWE}_{S, \sigma}\left(\dfrac{q}{\beta^i} M + E_i\right)  \Bigr \}_{i=1}^{l} \in \mathcal{R}_{\langle n, q \rangle }^{(k+1) \cdot l}$
\end{tcolorbox}

\begin{figure}[h!]
    \centering
  \includegraphics[width=1.0\linewidth]{figures/TFHE-fig2.pdf}
  \caption{GLev 密文结构示意图。}
  \label{fig:glev}
\end{figure}



\subsection{解密}

第一行的相位为 $(q/\beta)M+E_1\pmod q$。若 $\|E_1\|_\infty<q/(2\beta)$，最近整数解码逐系数恢复 $M\bmod\beta$。只有当原代表元能由其模 $\beta$ 的类唯一确定时，才能恢复它，例如系数属于 $\{0,\ldots,t-1\}$ 且 $t\le\beta$。类似地，若 $\|E_i\|_\infty<q/(2\beta^i)$，第 $i$ 行能解码模 $\beta^i$ 的结果。后面的行对误差的容忍度更低，但不一定无法解码；gadget 运算也不要求各行单独解码。

\subsection{Lev 与 RLev}

Lev 是 $n=1$ 时的 GLev，RLev 是 $k=1$ 时的 GLev。

